LLVManim is implemented as a three-layer pipeline---\textbf{Ingestion}, \textbf{Transformation}, and \textbf{Render}---orchestrated by a dedicated \textbf{CLI} entry point. The high-level structure is shown in Fig.~\ref{fig:architecture_detail}. Layer boundaries are enforced at build time by \texttt{import-linter}, ensuring that no module imports from a downstream layer. The entire codebase is written in Python~3.12 and managed with \texttt{uv}\cite{uv}.

\vspace{1ex}

\begin{figure*}[htbp]
\includegraphics[height=\textheight, width=\textwidth]{01_architecture_overview_vertical}
\caption{Detailed module-level architecture. The Ingestion layer collects LLVM IR and metadata, the Transformation layer maps it to an intermediate scene graph, and the Rendering layer renders animations or exports artifacts.}
\label{fig:architecture_detail}
\end{figure*}

\needspace{4\baselineskip}
\subsubsection{CLI Entry Point}
\label{sec:llvmanim:architecture:cli}
\subimport{architecture/}{cli.tex}

\vspace{1ex}

\needspace{4\baselineskip}
\subsubsection{Ingestion Layer}
\label{sec:llvmanim:architecture:ingestion}
\subimport{architecture/}{ingestion.tex}

\vspace{1ex}

\subsubsection{Transformation Layer}
\label{sec:llvmanim:architecture:transformation}
\subimport{architecture/}{transformation.tex}

\vspace{1ex}

\subsubsection{Render Layer}
\label{sec:llvmanim:architecture:render}
\subimport{architecture/}{render.tex}
